%% kapa: exam=1 round=33 year=2022 period=1 form=A subject=civil_law subject_ko=민법 seq=33 number=33
%% source: Q-Net 감정평가사 기출문제 (cst00309) artlSeq=5209997 file=raw/1008/5209997/2022년도 제33회 감정평가사 1차시험 1교시(민법,경제학,부동산학).pdf
%% conversion: kordoc 4.24.51 -> tools/gen1.py (auto); review=proofread-sonnet; flags=none
\begin{question}{33}
전세권에 관한 설명으로 옳지 않은 것은? (다툼이 있으면 판례에 따름)
\begin{choices}
  \item 타인의 토지에 있는 건물에 전세권을 설정한 때에는 전세권의 효력은 그 건물의 소유를 목적으로 한 지상권에 미친다.
  \item 건물전세권설정자가 건물의 존립을 위한 토지사용권을 가지지 못하여 그가 토지소유자의 건물철거 등 청구에 대항할 수 없는 경우, 전세권자는 토지소유자의 권리행사에 대항할 수 없다.
  \item 지상권을 가지는 건물소유자가 그 건물에 전세권을 설정하였으나 그가 2년 이상의 지료를 지급하지 아니하였음을 이유로 지상권설정자가 지상권의 소멸을 청구한 경우, 전세권자의 동의가 없다면 지상권은 소멸되지 않는다.
  \item 대지와 건물이 동일한 소유자에 속한 경우에 건물에 전세권을 설정한 때에는 그 대지 소유권의 특별승계인은 전세권설정자에 대하여 지상권을 설정한 것으로 본다.
  \item 건물에 대한 전세권의 존속기간을 1년 미만으로 정한 때에는 이를 1년으로 한다.
\end{choices}
\end{question}
\answer{3}
